% Work-MAT | 2. rocnik | hodina 15
\documentclass[11pt,a4paper]{article}
\input{../../sablony/pripravy/priprava.sty}
\usepackage{tikz}
\usetikzlibrary{arrows.meta}
\nastavPripravu{2. ročník}{Exponenciální rovnice}{15}{Exponenciální rovnice III: vytýkání}
\begin{document}
\poznamka{Metodu nejprve propoj s vytýkáním mnohočlenů: $t+2t=3t$. Rozdíl je v tom, že zde opakovanou veličinou je $a^x$. Orientačně 8 min motivace, 15 min vytýkání, 12 min procvičení, 10 min upevnění.}
\blok{Když se exponenciální výrazy sčítají}
Při součinu můžeme exponenty sčítat, ale při součtu \textbf{neplatí} $a^x+a^{x+1}=a^{2x+1}$. Místo toho užijeme $a^{x+1}=a\cdot a^x$ a výraz $a^x$ vytkneme.\par
\begin{prikladReseniBox}[title={Příklad 1}]
Vyřeš $2^{x+1}+2^x=24$.
\tcblower
\[\begin{aligned}2\cdot2^x+2^x&=24\\(2+1)2^x&=24\\2^x&=8\\2^x&=2^3\\x&=3.\end{aligned}\]
Zkouška: $2^4+2^3=16+8=24$, tedy $K=\{3\}$.
\end{prikladReseniBox}
\begin{prikladReseniBox}[title={Příklad 2}]
Vyřeš $3^{x+2}-3^x=72$.
\tcblower
\[\begin{aligned}9\cdot3^x-3^x&=72\\8\cdot3^x&=72\\3^x&=9\\x&=2.\end{aligned}\]
Zkouška: $3^4-3^2=81-9=72$.
\end{prikladReseniBox}
\begin{zadaniBox}[title={Věta - Vytknutí společné exponenciální mocniny}]
Pro $a>0$ lze $c\,a^{x+p}+d\,a^{x+q}$ upravit na $a^x(ca^p+da^q)$. Pokud je výraz v závorce nenulový, dostáváme rovnici pro jedinou mocninu $a^x$.
\end{zadaniBox}
\begin{prikladReseniBox}[title={Příklad 3}]
Vyřeš $2^{x+2}-2^{x-1}=28$.
\tcblower
\[\begin{aligned}(4-\tfrac12)2^x&=28\\\tfrac72\,2^x&=28\\2^x&=8\\x&=3.\end{aligned}\]
Zkouška: $2^5-2^2=32-4=28$.
\end{prikladReseniBox}
\poznamka{Při vytýkání někdy vystupuje koeficient záporný nebo nulový. S nulovým koeficientem by se celá exponenciální část vyrušila; nelze vždy bez kontroly dělit výrazem v závorce.}
\begin{zadaniReseniBox}[title={Samostatná práce 1}]
\begin{enumerate}[label=\alph*)]
\item $2^{x+3}+2^x=36$;
\item $5^{x+1}-5^x=100$;
\item $4^{x+1}+4^x=80$.
\end{enumerate}
\tcblower
\begin{enumerate}[label=\alph*)]
\item $(8+1)2^x=36$, tedy $2^x=4$, $x=2$.
\item $(5-1)5^x=100$, tedy $5^x=25$, $x=2$.
\item $(4+1)4^x=80$, tedy $4^x=16$, $x=2$.
\end{enumerate}
\end{zadaniReseniBox}
\begin{zadaniReseniBox}[title={Rozšíření 1}]
Vyřeš $3^{x+1}+3^{x-1}=30$.
\tcblower
$(3+\tfrac13)3^x=30$, tedy $\frac{10}{3}3^x=30$, odkud $3^x=9$ a $x=2$.
\end{zadaniReseniBox}
\begin{zadaniReseniBox}[title={Domácí úkol}]
\begin{enumerate}[label=\alph*)]
\item $2^{x+1}+3\cdot2^x=40$.
\item $7^{x+1}-6\cdot7^x=49$.
\item Uveď vlastní protipříklad k chybnému pravidlu $2^x+2^x=2^{2x}$.
\end{enumerate}
\tcblower
\begin{enumerate}[label=\alph*)]
\item $5\cdot2^x=40$, tedy $x=3$.
\item $7^x=49$, tedy $x=2$.
\item Například $x=2$: $4+4=8$, ale $2^4=16$.
\end{enumerate}
\end{zadaniReseniBox}
\end{document}
